\documentclass[11pt]{article}
\usepackage[a4paper, margin=2cm]{geometry}
\usepackage{graphicx}
\usepackage{amsmath}
\usepackage{enumitem}

% ============================================================================
%  HOW THIS FILE IS STRUCTURED — READ THIS BEFORE EDITING
%
%  There are TWO different "begin/end" pairs below and they are NOT the
%  same thing — mixing them up is the #1 cause of "no question blocks were
%  found" errors on upload.
%
%   1) \begin{document} ... \end{document}
%      Plain, standard LaTeX. It appears EXACTLY ONCE in the whole file —
%      right here, and again at the very bottom — and wraps everything.
%      Never add a second copy of it anywhere in the middle of the file.
%
%   2) \begin{question}{n} ... \end{question}
%      Our own custom tag (this is what the parser actually looks for).
%      It appears ONCE PER QUESTION, nested inside \begin{document} ...
%      \end{document}. Every question needs BOTH its own \begin{question}{n}
%      AND its own \end{question} placed right after that question's
%      content. The single \end{document} at the bottom of the file does
%      NOT close your questions for you — if even one question is missing
%      its \end{question}, the parser cannot find ANY question in the file
%      and you will see "no question blocks were found".
%
%  Rule of thumb: count your questions, then count your \end{question}
%  lines — they must match exactly.
% ============================================================================

\begin{document}

\subject{Physics}
\subjectcode{5054}
\paper{2}
\variant{1}
\session{M/J}
\year{2025}

\begin{question}{1}
\topic{Forces & Motion}
\marks{9}
\ref{5054/21/M/J/25 -- Q1}

\qtext
Fig. 1.1 shows a skydiver falling vertically through the air.

\image{q1-fig1.png}{Fig. 1.1}

In the first part of the fall, her speed increases and her acceleration decreases.

\part{2} On Fig. 1.2 sketch the speed--time graph for the skydiver.

\part{1} Explain how the graph shows that the acceleration decreases as the speed increases.
\endq
\end{question}
% ^^^ \end{question} above closes QUESTION 1. Every question needs this line.

\begin{question}{2}
\topic{Moments}
\marks{6}
\ref{5054/21/M/J/25 -- Q2}

\qtext
A uniform beam of length $2.0\ \text{m}$ is pivoted at its centre.

\part{2} State the principle of moments.

\part{4} Calculate the force needed at one end to balance a $5.0\ \text{N}$ weight placed $0.4\ \text{m}$ from the opposite end.
\endq
\end{question}
% ^^^ \end{question} above closes QUESTION 2.

\begin{question}{3}
\topic{Energy}
\marks{5}
\ref{5054/21/M/J/25 -- Q3}

\qtext
A ball of mass $0.20\ \text{kg}$ is dropped from a height of $1.8\ \text{m}$
above the ground. Take $g = 9.8\ \text{m/s}^2$ and ignore air resistance.

\part{2} Calculate the gravitational potential energy of the ball before it is released.

\part{3} Calculate the speed of the ball just before it hits the ground.
\endq
\end{question}
% ^^^ \end{question} above closes QUESTION 3.

% To add another question, copy one whole block above — from its opening
% \begin{question}{N} tag down to its own closing \end{question} tag — and
% edit the contents. (Written as {N} here on purpose: the parser scans the
% raw text and does not understand "%" comments, so writing a real digit
% like {4} in a comment would be read as an actual — and broken — question.)

\end{document}
% ^^^ This single \end{document} closes the FILE, not any individual
%     question — every \begin{question}{n} above must already have its own
%     \end{question} before this line.